\documentclass[12pt,a4paper]{article}

% ===== Encoding & Language =====
\usepackage[utf8]{inputenc}
\usepackage[T1]{fontenc}
\usepackage[czech]{babel}

% ===== Page Layout =====
\usepackage[a4paper, top=2.5cm, bottom=2.5cm, left=2.5cm, right=2.5cm]{geometry}

% ===== Fonts =====
\usepackage{mathptmx}
\usepackage[scaled=0.92]{helvet}

% ===== Spacing =====
\usepackage{setspace}
\onehalfspacing
\setlength{\parindent}{0pt}
\setlength{\parskip}{6pt}

% ===== Section Headings (sans-serif, bold) =====
\usepackage{titlesec}
\usepackage{needspace}
\titleformat{\section}{\Large\sffamily\bfseries}{\thesection}{1em}{}
\titleformat{\subsection}{\needspace{8\baselineskip}\large\sffamily\bfseries}{\thesubsection}{1em}{}
\titleformat{\subsubsection}{\needspace{5\baselineskip}\normalsize\sffamily\bfseries}{\thesubsubsection}{1em}{}
\titlespacing*{\section}{0pt}{18pt}{6pt}
\titlespacing*{\subsection}{0pt}{12pt}{6pt}
\titlespacing*{\subsubsection}{0pt}{12pt}{6pt}

% ===== Hyperlinks =====
\usepackage[unicode, colorlinks=true, linkcolor=black, urlcolor=blue, citecolor=black]{hyperref}

% ===== Graphics =====
\usepackage{graphicx}
\usepackage{float}

% ===== Code Listings =====
\usepackage{listings}
\usepackage{xcolor}

\definecolor{codebg}{RGB}{248,248,248}
\definecolor{codeframe}{RGB}{200,200,200}
\definecolor{codecomment}{RGB}{106,153,85}
\definecolor{codestring}{RGB}{163,21,21}
\definecolor{codekeyword}{RGB}{0,0,180}

\lstset{
    basicstyle=\small\ttfamily,
    backgroundcolor=\color{codebg},
    frame=single,
    rulecolor=\color{codeframe},
    breaklines=true,
    breakatwhitespace=false,
    tabsize=4,
    showstringspaces=false,
    numbers=none,
    aboveskip=10pt,
    belowskip=10pt,
    xleftmargin=2pt,
    xrightmargin=2pt,
    commentstyle=\color{codecomment},
    stringstyle=\color{codestring},
    keywordstyle=\color{codekeyword}\bfseries,
    literate=
        {á}{{\'a}}1 {é}{{\'e}}1 {í}{{\'{\i}}}1 {ó}{{\'o}}1 {ú}{{\'u}}1
        {ý}{{\'y}}1 {č}{{\v{c}}}1 {ď}{{\v{d}}}1 {ě}{{\v{e}}}1
        {ň}{{\v{n}}}1 {ř}{{\v{r}}}1 {š}{{\v{s}}}1 {ť}{{\v{t}}}1
        {ů}{{\r{u}}}1 {ž}{{\v{z}}}1
        {Á}{{\'A}}1 {É}{{\'E}}1 {Í}{{\'I}}1 {Ó}{{\'O}}1 {Ú}{{\'U}}1
        {Ý}{{\'Y}}1 {Č}{{\v{C}}}1 {Ď}{{\v{D}}}1 {Ě}{{\v{E}}}1
        {Ň}{{\v{N}}}1 {Ř}{{\v{R}}}1 {Š}{{\v{S}}}1 {Ť}{{\v{T}}}1
        {Ů}{{\r{U}}}1 {Ž}{{\v{Z}}}1
        {→}{{$\rightarrow$}}1,
}

% ===== Tables =====
\usepackage{tabularx}
\usepackage{booktabs}
\usepackage{array}

% ===== Lists =====
\usepackage{enumitem}
\setlist{nosep, leftmargin=2em}

% ===== Bibliography heading =====
\renewcommand{\refname}{Seznam použité literatury}

% ===== Widow/orphan prevention =====
\widowpenalty=10000
\clubpenalty=10000

% ========================================================================
\begin{document}
\raggedbottom

% ===== TITULNÍ STRANA =====
\begin{titlepage}
\begin{center}

\vspace*{0.5cm}
{\sffamily\Large\bfseries GYMNÁZIUM A STŘEDNÍ ODBORNÁ ŠKOLA}\\[4pt]
{\sffamily\large ul.\ Mládežníků 1115, Rokycany}

\vfill

{\sffamily\huge\bfseries Webová prezentace fotbalového klubu\\[6pt]
s~databází zápasů a~správou obsahu\\[6pt]
ve~WordPressu}

\vspace{1.5cm}

{\sffamily\Large maturitní práce}

\vfill

{\large
\begin{tabular}{@{}ll@{}}
Autor práce: & Lukáš Bejček \\
Vedoucí práce: & Mgr. Jaromír Háka \\
Obor: & Informační technologie \\
\end{tabular}
}

\vspace{1.5cm}

{\large Rokycany 2026}

\end{center}
\end{titlepage}

% ===== PROHLÁŠENÍ =====
\thispagestyle{empty}
\section*{Prohlášení}
\addcontentsline{toc}{section}{Prohlášení}

Prohlašuji, že jsem maturitní práci na téma \textit{Webová prezentace fotbalového klubu s~databází zápasů a~správou obsahu ve~WordPressu} vypracoval samostatně pod vedením vedoucího práce a~s~použitím odborné literatury a~dalších informačních zdrojů, které jsou citovány v~práci a~uvedeny v~seznamu použité literatury.

\vspace{2cm}

V~Mýtě dne \dotfill{} Podpis: \dotfill{}

\newpage

% ===== PODĚKOVÁNÍ =====
\thispagestyle{empty}
\section*{Poděkování}
\addcontentsline{toc}{section}{Poděkování}

Děkuji vedoucímu práce za odborné vedení, cenné rady a~připomínky při zpracování maturitní práce.

\newpage

% ===== ABSTRAKT =====
\thispagestyle{empty}
\section*{Abstrakt}
\addcontentsline{toc}{section}{Abstrakt}

Tato maturitní práce se zabývá návrhem a~implementací webové prezentace fotbalového klubu TJ Slavoj Mýto. Cílem bylo vytvořit moderní, responzivní web, který umožňuje správci klubu spravovat zápasy, týmy, hráče, galerie, sponzory a~kontakty bez znalosti programování. Web je postaven na redakčním systému WordPress s~vlastním tématem založeným na frameworku Bootstrap~5 a~vlastním pluginem pro rozšíření administrace. Místo externích placených pluginů byla zvolena vlastní implementace klíčových funkcí -- registrace typů obsahu, vlastních polí, filtrů a~uživatelských rolí -- přímo v~kódu tématu a~pluginu. Práce popisuje analýzu původního řešení, návrh datového modelu, implementaci šablon a~administračních nástrojů, testování, příručku administrátora a~postup nasazení na hosting.

\textbf{Klíčová slova:} WordPress, fotbalový klub, webová prezentace, custom post types, Bootstrap~5, PHP, responzivní design

\newpage

% ===== OBSAH =====
\thispagestyle{empty}
\tableofcontents
\newpage

% ========================================================================
\clearpage
\section{Úvod}

Fotbalový klub TJ Slavoj Mýto je amatérský sportovní klub působící v~obci Mýto u~Rokycan. Klub má několik týmů od mužů po mládežnické kategorie a~potřebuje moderní webovou prezentaci, na které budou přehledně zobrazeny zápasy, týmy, hráčské soupisky, fotogalerie, partneři klubu a~kontakty na vedení.

Původní web klubu byl charakteristický nepřehledným uspořádáním obsahu a~graficky zastaralým vzhledem. Cílem této práce je vytvořit graficky přijatelný a~přehledný web s~plně funkčním a~spravovatelným WordPressem, kde veškerý obsah lze spravovat přes standardní administrační rozhraní.

Práce je členěna do následujících kapitol. Kapitola~2 analyzuje původní řešení a~porovnává možné přístupy k~implementaci. Kapitola~3 popisuje implementaci -- datový model, šablony, administrační nástroje a~bezpečnost. Kapitola~4 dokumentuje testování a~ověření funkčnosti. Kapitola~5 slouží jako příručka administrátora pro správu obsahu webu. Kapitola~6 popisuje nasazení na produkční hosting. Kapitola~7 obsahuje závěr s~návrhy na rozšíření.

Celý projekt je dostupný jako open-source repozitář na GitHubu a~je připraven k~lokálnímu nasazení pomocí XAMPP.

% ========================================================================
\clearpage
\section{Analýza problému}

\subsection{Původní stav webu}

Původní web TJ Slavoj Mýto (\url{http://slavojmyto.cz/}) byl charakteristický nepřehledným uspořádáním obsahu a~graficky zastaralým vzhledem. Stránky obsahují sekce: úvodní stránku, historii klubu, výbor (kontakty), týmy se soupiskami, kalendář zápasů po sezónách, fotogalerie a~sponzory.

Ověřením živého webu byly zjištěny tyto nedostatky:

\begin{enumerate}
    \item \textbf{Nepřehledná navigace} -- struktura menu neodpovídá očekávání návštěvníka. Zápasy jsou skryté pod položkou \uv{Kalendář} (nikoliv \uv{Zápasy}), kontaktní údaje na vedení klubu pod \uv{Výbor} v~sekci \uv{O~klubu} (nikoliv \uv{Kontakty}). Chybí sekce aktualit/novinek.

    \item \textbf{Nepřehledný výpis zápasů} -- zápasy jsou zobrazeny jako dlouhý seznam seskupený podle týmů na jedné stránce za celou sezónu. Chybí jakékoliv filtrování (podle týmu, stavu zápasu) a~stránkování. Pro nalezení konkrétního zápasu je nutné procházet celou stránku.

    \item \textbf{Graficky zastaralý vzhled} -- vizuální styl webu neodpovídá moderním standardům webového designu. Layout se přizpůsobuje obrazovce pouze omezeně a~celkový grafický dojem působí zastarale.

    \item \textbf{Chybějící správa obsahu a~funkce} -- web postrádá standardní administrační rozhraní pro správu obsahu, vyhledávání, detaily jednotlivých zápasů, profily hráčů a~aktuality. Galerie používá dvouúrovňový systém (tým $\rightarrow$ ročník $\rightarrow$ album), ale bez filtrování.
\end{enumerate}

Z~původního webu bylo možné převzít vizuální identitu (logo, barvy klubu), strukturu sekcí a~dostupná data (soupisky hráčů, výsledky zápasů, kontakty na výbor, sponzory). Technická a~datová vrstva musela být navržena a~implementována od základu v~redakčním systému WordPress, který zajistí plnohodnotnou správu veškerého obsahu přes standardní administrační rozhraní.

\subsection{Požadavky na nové řešení}

Na základě zadání maturitní práce~\cite{zadani} a~potřeb klubu byly stanoveny tyto požadavky:

\begin{itemize}
    \item Evidence zápasů s~výsledky, střelci a~filtrováním podle sezóny, týmu a~stavu
    \item Evidence týmů se soupiskami hráčů a~realizačním týmem
    \item Fotogalerie členěná podle týmů a~sezón
    \item Kontakty na vedení klubu s~možností řazení
    \item Přehled sponzorů s~logy a~odkazy
    \item Stránka historie klubu
    \item Správa veškerého obsahu přes administraci WordPress bez zásahu do kódu
    \item Responzivní design (desktop i~mobilní zařízení)
    \item Minimální JavaScript -- preferovat serverové řešení v~PHP
\end{itemize}

\subsection{Volba technologií}

\subsubsection{CMS -- WordPress}

Pro implementaci byl zvolen redakční systém WordPress~\cite{wordpress}, nejrozšířenější CMS na světě s~podílem přes 40~\% všech webových stránek. WordPress nabízí:

\begin{itemize}
    \item Vestavěné administrační rozhraní pro správu obsahu
    \item Systém custom post types a~taxonomií pro strukturovaná data
    \item Šablonovou hierarchii pro mapování URL na šablony
    \item REST API pro případné budoucí rozšíření
    \item Rozsáhlou dokumentaci a~komunitu~\cite{wpdev}
\end{itemize}

\subsubsection{CSS framework -- Bootstrap~5}

Pro responzivní layout byl zvolen framework Bootstrap~5~\cite{bootstrap}, který poskytuje:

\begin{itemize}
    \item Grid systém pro responzivní rozvržení stránky
    \item Utility třídy pro rychlé stylování (spacing, barvy, typography)
    \item Hotové komponenty (navbar, cards, badges, forms)
    \item Mobile-first přístup k~responzivnímu designu
\end{itemize}

Bootstrap je načítán z~CDN~\cite{jsdelivr}, čímž odpadá nutnost build nástrojů a~zároveň se využívá cache prohlížeče.

\subsubsection{Lokální vývojové prostředí -- XAMPP}

Pro lokální vývoj a~testování byl zvolen balík XAMPP~\cite{xampp}, který obsahuje Apache webserver, MySQL databázi a~PHP interpret v~jedné instalaci. XAMPP je dostupný zdarma pro Windows, macOS i~Linux.

\subsubsection{Grafický návrh -- Figma}

Grafický návrh webu byl vytvořen v~nástroji Figma~\cite{figma} ve variantách pro desktop i~mobilní zařízení.

\subsection{Vlastní implementace vs.\ externí pluginy}

Zadání práce doporučuje řadu externích pluginů: Custom Post Type UI~\cite{cptui} pro registraci CPT, Advanced Custom Fields~\cite{acf} pro správu vlastních polí, FacetWP pro filtrování obsahu, The Events Calendar pro kalendář zápasů, Modula pro galerie a~User Role Editor pro správu rolí.

Po analýze jsem zvolili \textbf{vlastní implementaci} klíčových funkcí přímo v~kódu tématu a~vlastním pluginu. Důvody:

\begin{table}[H]
\centering
\begin{tabularx}{\textwidth}{l X X}
\toprule
\textbf{Aspekt} & \textbf{Externí pluginy} & \textbf{Vlastní implementace} \\
\midrule
Kontrola nad kódem & Omezená -- závislost na cizím kódu & Plná -- kód je čitelný a~obhajitelný \\
Závislosti & Více pluginů = více bodů selhání & Pouze WordPress core + vlastní kód \\
Výkon & Každý plugin přidává zátěž & Minimální overhead \\
Údržba & Nutnost aktualizovat pluginy & Kód je pod vlastní správou \\
Licence & Některé pluginy jsou placené & Vše zdarma (GPL-2.0) \\
\bottomrule
\end{tabularx}
\caption{Srovnání externích pluginů a~vlastní implementace}
\end{table}

Jediným nainstalovaným pluginem je \textbf{Slavoj Custom Fields} -- vlastní plugin vytvořený pro tento projekt, který rozšiřuje administrační prostředí.

\subsection{Custom post types vs.\ vlastní databázové tabulky}

WordPress nabízí dva přístupy k~ukládání strukturovaných dat: custom post types (CPT) s~meta poli, nebo vlastní databázové tabulky.

Pro tento projekt byly zvoleny CPT z~následujících důvodů~\cite{wpdev}:

\begin{enumerate}
    \item \textbf{Integrace s~WordPress admin} -- CPT automaticky získají administrační rozhraní (přehled, editace, mazání, koš) bez nutnosti psát vlastní CRUD logiku.
    \item \textbf{Podpora taxonomií} -- termy (sezóna, kategorie týmu) lze přiřazovat nativně.
    \item \textbf{WP\_Query} -- dotazování probíhá standardním WordPress API, které řeší caching, stránkování i~bezpečnost.
    \item \textbf{REST API} -- parametr \texttt{show\_in\_rest => true} zpřístupní data přes JSON API.
    \item \textbf{Šablonová hierarchie} -- WordPress automaticky mapuje URL na šablony (\texttt{/zapasy/} $\rightarrow$ \texttt{archive-zapas.php}).
\end{enumerate}

Alternativa s~vlastními tabulkami by vyžadovala ruční implementaci administrace, validace, stránkování a~URL routingu.

\subsection{Archivní šablony vs.\ stránkové šablony}

Pro zobrazení seznamů CPT (zápasy, týmy, galerie) byly zvoleny \textbf{archivní šablony} (\texttt{archive-zapas.php}) místo stránkových šablon (\texttt{page-zapasy.php}). Důvody~\cite{wpdev}:

\begin{itemize}
    \item WordPress automaticky vytvoří čistou URL \texttt{/zapasy/} bez nutnosti zakládat stránku v~administraci
    \item Stránkování funguje nativně
    \item Taxonomické filtry se integrují přirozeně přes GET parametry
\end{itemize}

Stránkové šablony zůstávají pouze pro statický obsah: historie, kontakty, sponzoři.

% ========================================================================
\clearpage
\section{Implementace}

\subsection{Architektura řešení}

Projekt se skládá ze dvou komponent:

\textbf{Téma \texttt{tj-slavoj-myto}} obsahuje prezentační logiku a~registraci datového modelu:

\begin{itemize}
    \item registrace custom post types a~taxonomií
    \item šablony pro front-end (homepage, archivy, detaily)
    \item CSS styly rozdělené do komponentových souborů
    \item pomocné funkce (řazení, vazby, počítání hráčů)
    \item globální filtry pro konzistentní chování taxonomií
\end{itemize}

\textbf{Plugin \texttt{slavoj-custom-fields}} poskytuje administrační nástroje:

\begin{itemize}
    \item vlastní sloupce v~admin přehledech (datum, tým, skóre, číslo dresu)
    \item dropdown filtry pro filtrování v~administraci
    \item řaditelné sloupce (datum zápasu, číslo dresu)
    \item nástrojová stránka pro správu a~seedování ukázkových dat
\end{itemize}

Toto rozdělení odpovídá doporučené WordPress architektuře~\cite{wpdev}: téma řeší \uv{jak vypadá}, plugin řeší \uv{jak se spravuje}.

\subsection{Adresářová struktura}

\begin{lstlisting}
web/
+-- wp-content/
    |-- themes/tj-slavoj-myto/
    |   |-- functions.php              # Registrace CPT, taxonomii, helperu
    |   |-- front-page.php             # Homepage
    |   |-- archive-zapas.php          # Seznam zapasu s filtry
    |   |-- archive-tym.php            # Seznam tymu
    |   |-- archive-galerie.php        # Prehled fotoalb
    |   |-- single-zapas.php           # Detail zapasu
    |   |-- single-tym.php             # Profil tymu se soupiskou
    |   |-- single-hrac.php            # Profil hrace
    |   |-- single-galerie.php         # Fotoalbum s lightboxem
    |   |-- page-historie.php          # Historie klubu
    |   |-- page-kontakty.php          # Kontakty vyboru
    |   |-- page-sponzori.php          # Partneri klubu
    |   |-- template-parts/            # Znovupouzitelne komponenty
    |   +-- assets/css/                # Stylove soubory
    +-- plugins/slavoj-custom-fields/
        +-- slavoj-custom-fields.php   # Administracni plugin
\end{lstlisting}

\subsection{Datový model}

\subsubsection{Custom post types}

Projekt definuje šest vlastních typů obsahu registrovaných v~\texttt{functions.php} pomocí funkce \texttt{register\_post\_type()}~\cite{wpdev}:

\begin{table}[H]
\centering
\small
\begin{tabularx}{\textwidth}{l l X l}
\toprule
\textbf{CPT} & \textbf{Slug} & \textbf{Meta pole} & \textbf{Účel} \\
\midrule
Zápasy & \texttt{zapas} & \texttt{datum\_zapasu}, \texttt{cas\_zapasu}, \texttt{domaci}, \texttt{hoste}, \texttt{skore}, \texttt{strelci} & Evidence zápasů \\
Týmy & \texttt{tym} & \texttt{hlavni\_trener}, \texttt{asistent\_trenera}, \texttt{zdravotnik}, \texttt{tym\_slug} & Profily týmů \\
Hráči & \texttt{hrac} & \texttt{cislo}, \texttt{rok\_narozeni}, \texttt{tym\_slug} & Soupisky hráčů \\
Galerie & \texttt{galerie} & --- & Fotoalba \\
Sponzoři & \texttt{sponzor} & \texttt{web\_sponzora} & Partneři klubu \\
Kontakty & \texttt{kontakt} & \texttt{pozice}, \texttt{telefon}, \texttt{email}, \texttt{poradi} & Výbor klubu \\
\bottomrule
\end{tabularx}
\caption{Přehled custom post types}
\end{table}

Ukázka registrace CPT pro zápasy:

\begin{lstlisting}[language=PHP]
register_post_type('zapas', array(
    'labels' => array(
        'name'               => 'Zápasy',
        'singular_name'      => 'Zápas',
        'add_new'            => 'Přidat zápas',
        'add_new_item'       => 'Přidat nový zápas',
        // ...
    ),
    'public'              => true,
    'has_archive'         => true,
    'rewrite'             => array('slug' => 'zapasy'),
    'supports'            => array('title', 'thumbnail'),
    'menu_icon'           => 'dashicons-awards',
    'show_in_rest'        => true,
    'taxonomies'          => array('sezona', 'kategorie-tymu', 'stav-zapasu'),
    'capability_type'     => array('zapas', 'zapasy'),
    'map_meta_cap'        => true,
));
\end{lstlisting}

Každý CPT má vlastní \texttt{capability\_type} s~\texttt{map\_meta\_cap => true}, což umožňuje granulární řízení oprávnění pro jednotlivé uživatelské role.

\subsubsection{Taxonomie}

Projekt používá čtyři vlastní taxonomie registrované přes \texttt{register\_taxonomy()}~\cite{wpdev}:

\begin{table}[H]
\centering
\small
\begin{tabularx}{\textwidth}{l l l X}
\toprule
\textbf{Taxonomie} & \textbf{Slug} & \textbf{Sdílená mezi CPT} & \textbf{Účel} \\
\midrule
Sezóna & \texttt{sezona} & zapas, tym, hrac, galerie & Filtrování podle ročníku \\
Kategorie týmu & \texttt{kategorie-tymu} & zapas, tym, hrac, galerie & Muži A, B, Dorost, Žáci\ldots \\
Stav zápasu & \texttt{stav-zapasu} & zapas & Nadcházející, Odehraný, Zrušený \\
Pozice hráče & \texttt{pozice-hrace} & hrac & Brankáři, Hráči v~poli \\
\bottomrule
\end{tabularx}
\caption{Přehled taxonomií}
\end{table}

Ukázka registrace taxonomie:

\begin{lstlisting}[language=PHP]
register_taxonomy('sezona', array('zapas', 'tym', 'hrac', 'galerie'), array(
    'labels'            => array(
        'name'          => 'Sezóny',
        'singular_name' => 'Sezóna',
    ),
    'hierarchical'      => false,
    'public'            => true,
    'show_admin_column' => true,
    'show_in_rest'      => true,
));
\end{lstlisting}

Sdílení taxonomií mezi více CPT umožňuje jednotnou klasifikaci obsahu a~zjednodušuje implementaci filtrování i~návazných výpisů napříč jednotlivými sekcemi webu.

\subsubsection{Vazby mezi entitami}

Vztahy mezi typy obsahu jsou řešeny dvěma mechanismy:

\textbf{Taxonomická vazba (M:N)} -- sezóna a~kategorie týmu jsou sdíleny mezi zápasy, týmy, hráči a~galeriemi pomocí WordPress taxonomií.

\textbf{Meta pole vazba (1:N)} -- hráči jsou propojeni s~týmy přes meta pole \texttt{tym\_slug}. Každý hráč má uložen slug svého týmu (např.\ \texttt{muzi-a}), který odpovídá hodnotě \texttt{tym\_slug} u~příslušného záznamu CPT \texttt{tym}.

\begin{lstlisting}
Zapas --(taxonomie)--> Sezona <--(taxonomie)-- Tym
  |                                              |
  +--(taxonomie)--> Kategorie tymu <--(taxonomie)+
                                                 |
                                           (meta: tym_slug)
                                                 |
                                               Hrac
\end{lstlisting}

Funkce \texttt{slavoj\_count\_hracu\_tymu()} automaticky počítá hráče přiřazené k~týmu a~výsledek ukládá do meta pole \texttt{pocet\_hracu}.

\subsubsection{ER diagram}

\begin{lstlisting}
+--------------------------------+
|            SEZONA              |  taxonomie: sezona
|  slug: 2024-25 / 2025-26 ...  |
+----------+---------------------+
           | m:n (term relationship)
     +-----+------+-----------------+
     |     |      |                 |
     v     v      v                 v
+--------+ +--------+ +----------+ +---------+
| ZAPAS  | |  TYM   | |  HRAC    | | GALERIE |
+---+----+ +---+----+ +----+-----+ +---------+
    |          |            |
    |          |  1:n       | n:1 (tym_slug)
    |          +------------+
    v m:n
+--------------------------------+
|         KATEGORIE TYMU         |  taxonomie: kategorie-tymu
+--------------------------------+

+---------+   +----------+
| KONTAKT |   | SPONZOR  |   (samostatne entity)
+---------+   +----------+
\end{lstlisting}

\subsection{Vlastní meta pole (nahrazuje Advanced Custom Fields)}

Místo pluginu Advanced Custom Fields~\cite{acf} jsou meta boxy implementovány přímo v~PHP pomocí WordPress funkcí \texttt{add\_meta\_box()}, \texttt{get\_post\_meta()} a~\texttt{update\_post\_meta()}~\cite{wpdev}.

Ukázka definice meta boxu pro zápasy:

\begin{lstlisting}[language=PHP]
add_meta_box(
    'slavoj_zapas_detail',
    'Detail zápasu',
    'slavoj_zapas_meta_box_html',
    'zapas',
    'normal',
    'high'
);
\end{lstlisting}

Pole pro jednotlivé CPT:

\textbf{Zápas:} datum zápasu (date), čas výkopu (time), domácí tým (text), hostující tým (text), skóre (text), střelci (text).

\textbf{Tým:} slug týmu (text), hlavní trenér (text), asistent trenéra (text), zdravotník (text), počet hráčů (automaticky vypočteno).

\textbf{Hráč:} číslo dresu (number 1--99), rok narození (number), slug týmu (text).

\textbf{Kontakt:} funkce/pozice (text), telefon (text), e-mail (email), pořadí zobrazení (number).

\textbf{Sponzor:} webové stránky (url).

\subsection{Filtrace obsahu}

Místo pluginů FacetWP nebo Search \& Filter Pro~\cite{facetwp} jsou filtry implementovány jako HTML formuláře s~GET parametry:

\begin{lstlisting}[language=HTML]
<form method="get"
      action="<?php echo esc_url(get_post_type_archive_link('zapas')); ?>">
    <select name="kat" onchange="this.form.submit()">
        <option value="">Všechny týmy</option>
        <!-- dynamicky generované možnosti -->
    </select>
    <noscript><button type="submit">Filtrovat</button></noscript>
</form>
\end{lstlisting}

PHP kód v~šabloně načte GET parametry, sanitizuje je pomocí \texttt{sanitize\_text\_field()} a~předá do \texttt{WP\_Query} s~\texttt{tax\_query} a~\texttt{meta\_query}~\cite{wpdev}.

Šablony s~filtrací:

\begin{table}[H]
\centering
\begin{tabular}{ll}
\toprule
\textbf{Šablona} & \textbf{Filtry} \\
\midrule
\texttt{archive-zapas.php} & Tým, sezóna, stav zápasu \\
\texttt{archive-tym.php} & Sezóna \\
\texttt{archive-galerie.php} & Tým, sezóna \\
\bottomrule
\end{tabular}
\caption{Šablony s~filtrací obsahu}
\end{table}

Formuláře používají \texttt{onchange="this.form.submit()"} pro automatické odeslání při změně výběru, s~\texttt{<noscript>} fallbackem pro prohlížeče bez JavaScriptu.

\subsection{Řazení a~kontextové filtrování taxonomií}

\subsubsection{Problém}

WordPress řadí termy taxonomií abecedně. Pro fotbalový klub je přirozené hierarchické pořadí: Muži~A $\rightarrow$ Muži~B $\rightarrow$ Dorost $\rightarrow$ Starší žáci $\rightarrow$ Mladší žáci $\rightarrow$ Přípravky. Navíc kategorie \uv{Stará garda} se má zobrazovat pouze u~galerií, nikoli u~zápasů nebo hráčů.

\subsubsection{Řešení}

Pořadí je definováno centrálně ve funkci \texttt{slavoj\_kategorie\_poradi()}:

\begin{lstlisting}[language=PHP]
function slavoj_kategorie_poradi() {
    return array(
        'muzi-a'            => 'Muži A',
        'muzi-b'            => 'Muži B',
        'stara-garda'       => 'Stará garda',
        'dorost'            => 'Dorost',
        'starsi-zaci'       => 'Starší žáci',
        'mladsi-zaci'       => 'Mladší žáci',
        'starsi-pripravka'  => 'Starší přípravka',
        'mladsi-pripravka'  => 'Mladší přípravka',
        'mini-pripravka'    => 'Mini přípravka',
    );
}
\end{lstlisting}

Globální filtr na WordPress hook \texttt{get\_terms} automaticky seřadí termy podle tohoto kanonického pořadí a~odfiltruje kontextově omezené termy (např.\ \uv{Stará garda} se zobrazí pouze v~kontextu galerií).

\subsection{Šablony}

\subsubsection{Přehled šablon}

\begin{table}[H]
\centering
\small
\begin{tabularx}{\textwidth}{l l X}
\toprule
\textbf{Šablona} & \textbf{URL} & \textbf{Obsah} \\
\midrule
\texttt{front-page.php} & \texttt{/} & Nadcházející zápasy, aktuality \\
\texttt{archive-zapas.php} & \texttt{/zapasy/} & Filtrovaný seznam zápasů \\
\texttt{archive-tym.php} & \texttt{/tymy/} & Karty týmů \\
\texttt{archive-galerie.php} & \texttt{/galerie/} & Grid fotoalb \\
\texttt{single-zapas.php} & \texttt{/zapasy/\{slug\}/} & Detail zápasu \\
\texttt{single-tym.php} & \texttt{/tymy/\{slug\}/} & Profil týmu, soupiska \\
\texttt{single-hrac.php} & \texttt{/hraci/\{slug\}/} & Profil hráče \\
\texttt{single-galerie.php} & \texttt{/galerie/\{slug\}/} & Fotoalbum s~lightboxem \\
\texttt{page-historie.php} & \texttt{/historie/} & Historie klubu \\
\texttt{page-kontakty.php} & \texttt{/kontakty/} & Kontakty výboru \\
\texttt{page-sponzori.php} & \texttt{/sponzori/} & Loga partnerů \\
\bottomrule
\end{tabularx}
\caption{Přehled šablon}
\end{table}

\subsubsection{Znovupoužitelné komponenty (template-parts)}

\begin{itemize}
    \item \textbf{\texttt{card-match.php}} -- karta zápasu (domácí vs hosté, datum, skóre, střelci)
    \item \textbf{\texttt{hero-team.php}} -- hero sekce s~názvem týmu a~trenéry
    \item \textbf{\texttt{site-header.php}} -- hlavička webu s~navigací
    \item \textbf{\texttt{site-footer.php}} -- patička webu
\end{itemize}

\subsection{CSS architektura}

Styly jsou rozdělené do komponentových souborů a~načítány přes \texttt{wp\_enqueue\_style()}~\cite{wpdev} s~definovanými závislostmi:

\begin{lstlisting}[language=PHP]
function slavoj_enqueue_scripts() {
    $ver = wp_get_theme()->get('Version');
    wp_enqueue_style('bootstrap',
        'https://cdn.jsdelivr.net/npm/bootstrap@5.3.3/dist/css/bootstrap.min.css',
        array(), '5.3.3');
    wp_enqueue_script('bootstrap-js',
        'https://cdn.jsdelivr.net/npm/bootstrap@5.3.3/dist/js/bootstrap.bundle.min.js',
        array(), '5.3.3', true);

    $css = get_template_directory_uri() . '/assets/css/';
    wp_enqueue_style('slavoj-utilities',  $css . 'utilities.css',
        array('bootstrap'), $ver);
    wp_enqueue_style('slavoj-header', $css . 'components/header.css',
        array('slavoj-utilities'), $ver);
    wp_enqueue_style('slavoj-buttons', $css . 'components/buttons.css',
        array('slavoj-utilities'), $ver);
    wp_enqueue_style('slavoj-cards', $css . 'components/cards.css',
        array('slavoj-utilities'), $ver);
    wp_enqueue_style('slavoj-main', $css . 'main.css',
        array('slavoj-utilities', 'slavoj-buttons',
              'slavoj-cards', 'slavoj-header'), $ver);
}
add_action('wp_enqueue_scripts', 'slavoj_enqueue_scripts');
\end{lstlisting}

Soubor \texttt{style.css} v~kořeni tématu slouží pouze jako identifikátor pro WordPress (název tématu, verze, autor) a~neobsahuje žádné styly. Veškeré styly a~skripty jsou načítány přes \texttt{wp\_enqueue\_style()} a~\texttt{wp\_enqueue\_script()} -- nikdy přímo v~šablonách~\cite{wpdev}.

\begin{lstlisting}
assets/css/
|-- utilities.css              # Barvy, spacing, utility tridy
|-- main.css                   # Hlavni styly sablon
+-- components/
    |-- header.css             # Navigace a hlavicka
    |-- nav-mobile.css         # Mobilni menu
    |-- buttons.css            # Tlacitka
    |-- hero.css               # Hero sekce
    +-- cards.css              # Karty zapasu, galerii, hracu
\end{lstlisting}

\subsection{Plugin Slavoj Custom Fields}

Plugin je uložen v~souboru \texttt{slavoj-custom-fields.php} a~poskytuje administrační nástroje:

\textbf{Vlastní sloupce v~přehledech:}
\begin{itemize}
    \item Zápasy: datum (řaditelný chronologicky), tým, skóre
    \item Hráči: číslo dresu (řaditelné numericky), pozice, tým
\end{itemize}

Ukázka implementace vlastního sloupce:

\begin{lstlisting}[language=PHP]
function slavoj_zapas_admin_column_data($column, $post_id) {
    if ($column === 'datum_zapasu') {
        $datum = get_post_meta($post_id, 'datum_zapasu', true);
        if ($datum) {
            $dt = DateTime::createFromFormat('Y-m-d', $datum);
            echo $dt ? esc_html($dt->format('j. n. Y'))
                     : esc_html($datum);
        }
    }
    // ...
}
\end{lstlisting}

\textbf{Dropdown filtry:} nad seznamy příspěvků lze filtrovat podle sezóny, kategorie týmu nebo týmu (u~hráčů).

\textbf{Nástrojová stránka} (Nástroje $\rightarrow$ Slavoj nastavení):
\begin{itemize}
    \item přehled registrovaných CPT a~taxonomií
    \item tlačítko \uv{Vytvořit výchozí hodnoty taxonomií}
    \item tlačítko \uv{Vytvořit ukázková data} (zápasy, hráči, kontakty, galerie)
    \item tlačítko \uv{Vytvořit stránky webu}
\end{itemize}

\subsection{Uživatelské role}

Místo pluginu User Role Editor~\cite{roleeditor} je vlastní role registrována přímo v~kódu:

\begin{lstlisting}[language=PHP]
add_role('slavoj_editor', 'Správce obsahu', array(
    'read'                   => true,
    'edit_posts'             => true,
    'edit_others_posts'      => true,
    'publish_posts'          => true,
    'upload_files'           => true,
    'manage_categories'      => true,
    // + oprávnění pro všechny CPT
));
\end{lstlisting}

\begin{table}[H]
\centering
\begin{tabularx}{\textwidth}{l l X}
\toprule
\textbf{Role} & \textbf{Popis} & \textbf{Oprávnění} \\
\midrule
Administrátor & Správce webu & Plný přístup \\
Správce obsahu & Správce zápasů a~obsahu & Editace CPT, nahrávání médií, bez nastavení WP \\
\bottomrule
\end{tabularx}
\caption{Uživatelské role}
\end{table}

\subsection{Optimalizace výkonu}

Základní optimalizace jsou implementovány přímo v~kódu tématu:

\begin{table}[H]
\centering
\begin{tabularx}{\textwidth}{l X}
\toprule
\textbf{Optimalizace} & \textbf{Implementace} \\
\midrule
Lazy loading obrázků & Funkce \texttt{slavoj\_add\_lazy\_loading()} přidává \texttt{loading="lazy"}~\cite{lazyload} \\
Bootstrap z~CDN & CSS a~JS z~jsdelivr.net -- využívá cache prohlížeče~\cite{jsdelivr} \\
Minimální JavaScript & Filtry přes HTML formuláře, bez AJAX knihoven \\
Mobile-first CSS & Vlastní CSS psaný mobile-first, minimální velikost \\
\bottomrule
\end{tabularx}
\caption{Optimalizace výkonu}
\end{table}

\begin{lstlisting}[language=PHP]
function slavoj_add_lazy_loading($content) {
    return str_replace('<img ', '<img loading="lazy" ', $content);
}
add_filter('the_content', 'slavoj_add_lazy_loading');
add_filter('post_thumbnail_html', 'slavoj_add_lazy_loading');
\end{lstlisting}

\subsection{Bezpečnost}

Projekt dodržuje standardní bezpečnostní praktiky WordPressu~\cite{wpdev,themehandbook}:

\begin{itemize}
    \item \textbf{Sanitizace vstupů:} Všechna data z~formulářů procházejí přes \texttt{sanitize\_text\_field()} nebo \texttt{sanitize\_email()}.
    \item \textbf{Escapování výstupů:} Veškerý výstup do HTML používá \texttt{esc\_html()}, \texttt{esc\_attr()} nebo \texttt{esc\_url()}.
    \item \textbf{Nonce ochrana:} Každý meta box formulář je chráněn WordPress nonce tokenem (\texttt{wp\_nonce\_field} / \texttt{wp\_verify\_nonce}).
    \item \textbf{WP\_Query:} Databázové dotazy využívají výhradně WordPress API, nikoli přímé SQL.
    \item \textbf{Capability mapping:} Každý CPT má vlastní \texttt{capability\_type} pro granulární řízení oprávnění.
    \item \textbf{WPS Hide Login:} Bezplatný plugin z~oficiálního repozitáře WordPressu, který mění výchozí přihlašovací URL \texttt{/wp-login.php} na vlastní adresu (např.\ \texttt{/prihlaseni/}). Přístup na \texttt{/wp-login.php} i~\texttt{/wp-admin/} (pro nepřihlášené) vrací chybu~404. Tím se výrazně snižuje riziko automatizovaných brute-force útoků, protože útočníci nemohou najít standardní přihlašovací stránku. Nastavení pluginu: \textbf{Nastavení $\rightarrow$ WPS Hide Login} -- stačí zadat požadovaný slug a~uložit.
\end{itemize}

\subsection{JavaScript}

Projekt záměrně minimalizuje použití JavaScriptu. Celkový rozsah vlastního JS je přibližně 70~řádků~\cite{jsreport}:

\begin{enumerate}
    \item \textbf{Filtry} -- \texttt{onchange="this.form.submit()"} s~\texttt{<noscript>} fallbackem (1~řádek na formulář)
    \item \textbf{Lightbox galerie} v~\texttt{single-galerie.php} (cca 55~řádků) -- kliknutí na foto, navigace šipkami/Escape
    \item \textbf{Quick-add střelců} v~\texttt{single-zapas.php} (cca 12~řádků) -- administrační pomůcka
\end{enumerate}

Všechny funkce mají funkční fallback bez JavaScriptu. Žádné frameworky ani build nástroje nejsou použity.

% ========================================================================
\clearpage
\section{Testování}

\subsection{Přehled testování}

Testování probíhalo průběžně během vývoje. Byly provedeny následující typy testů:

\begin{enumerate}
    \item \textbf{Funkční testování} -- ověření správné funkčnosti všech stránek, filtrů a~administrace
    \item \textbf{Bezpečnostní kontrola kódu} -- kontrola sanitizace, escapování a~nonce ochrany
    \item \textbf{Responzivní testování} -- ověření zobrazení na různých zařízeních
    \item \textbf{Kontrola WordPress best practices} -- dodržení standardů pro vývoj témat~\cite{themehandbook}
\end{enumerate}

\subsection{Funkční testovací scénáře}

\subsubsection*{Scénář 1: Přidání nového zápasu}

\begin{table}[H]
\centering
\small
\begin{tabularx}{\textwidth}{c X X c}
\toprule
\textbf{Krok} & \textbf{Akce} & \textbf{Očekávaný výsledek} & \textbf{Stav} \\
\midrule
1 & Přihlášení do administrace & Zobrazí se dashboard & Ano \\
2 & Klik na Zápasy $\rightarrow$ Přidat zápas & Zobrazí se formulář & Ano \\
3 & Vyplnění polí (datum, čas, týmy, skóre) & Pole přijímají data & Ano \\
4 & Přiřazení sezóny a~kategorie týmu & Taxonomie přiřazeny & Ano \\
5 & Klik na Publikovat & Zápas uložen & Ano \\
6 & Zobrazení na \texttt{/zapasy/} & Zápas se zobrazí v~seznamu & Ano \\
7 & Klik na detail zápasu & Detail s~daty se zobrazí & Ano \\
\bottomrule
\end{tabularx}
\end{table}

\subsubsection*{Scénář 2: Filtrování zápasů}

\begin{table}[H]
\centering
\small
\begin{tabularx}{\textwidth}{c X X c}
\toprule
\textbf{Krok} & \textbf{Akce} & \textbf{Očekávaný výsledek} & \textbf{Stav} \\
\midrule
1 & Otevření \texttt{/zapasy/} & Zobrazí se všechny zápasy & Ano \\
2 & Výběr \uv{Muži A} ve filtru týmu & Stránka se přenačte, zobrazí pouze Muži~A & Ano \\
3 & Výběr \uv{2025/26} ve filtru sezóny & Zobrazí pouze zápasy dané sezóny & Ano \\
4 & Výběr \uv{Odehrané} ve filtru stavu & Zobrazí pouze odehrané zápasy & Ano \\
5 & Reset filtrů (výběr \uv{Všechny}) & Zobrazí se všechny zápasy & Ano \\
\bottomrule
\end{tabularx}
\end{table}

\subsubsection*{Scénář 3: Soupiska týmu}

\begin{table}[H]
\centering
\small
\begin{tabularx}{\textwidth}{c X X c}
\toprule
\textbf{Krok} & \textbf{Akce} & \textbf{Očekávaný výsledek} & \textbf{Stav} \\
\midrule
1 & Otevření \texttt{/tymy/} & Zobrazí se karty týmů & Ano \\
2 & Klik na tým & Detail týmu se soupiskou & Ano \\
3 & Ověření počtu hráčů & Odpovídá počtu hráčů v~databázi & Ano \\
4 & Kontrola řazení hráčů & Seřazeni podle čísla dresu vzestupně & Ano \\
\bottomrule
\end{tabularx}
\end{table}

\subsubsection*{Scénář 4: Galerie}

\begin{table}[H]
\centering
\small
\begin{tabularx}{\textwidth}{c X X c}
\toprule
\textbf{Krok} & \textbf{Akce} & \textbf{Očekávaný výsledek} & \textbf{Stav} \\
\midrule
1 & Otevření \texttt{/galerie/} & Zobrazí se přehled alb & Ano \\
2 & Klik na album & Detail s~fotogalerií & Ano \\
3 & Klik na fotografii & Lightbox se otevře & Ano \\
4 & Navigace šipkami/Escape & Funkční navigace, zavření & Ano \\
\bottomrule
\end{tabularx}
\end{table}

\subsubsection*{Scénář 5: Uživatelská role -- Správce obsahu}

\begin{table}[H]
\centering
\small
\begin{tabularx}{\textwidth}{c X X c}
\toprule
\textbf{Krok} & \textbf{Akce} & \textbf{Očekávaný výsledek} & \textbf{Stav} \\
\midrule
1 & Přihlášení jako Správce obsahu & Dashboard bez nastavení WP & Ano \\
2 & Editace zápasu & Povoleno & Ano \\
3 & Nahrání obrázku & Povoleno & Ano \\
4 & Přístup k~Nastavení $\rightarrow$ Obecné & Zamítnuto & Ano \\
5 & Přístup k~Pluginy & Zamítnuto & Ano \\
\bottomrule
\end{tabularx}
\end{table}

\subsection{Bezpečnostní kontrola}

Kontrola kódu provedená v~březnu 2026 potvrdila dodržení bezpečnostních standardů~\cite{themehandbook}:

\begin{itemize}
    \item Sanitizace vstupů: \texttt{sanitize\_text\_field()}, \texttt{sanitize\_email()}, \texttt{wp\_unslash()}
    \item Escapování výstupů: \texttt{esc\_html()}, \texttt{esc\_attr()}, \texttt{esc\_url()}
    \item Nonce ochrana formulářů: \texttt{wp\_nonce\_field()}, \texttt{wp\_verify\_nonce()}
    \item \texttt{WP\_Query} místo raw SQL
    \item Správné použití \texttt{wp\_reset\_postdata()} po custom queries
\end{itemize}

\subsection{Nalezené a~opravené chyby}

Během testování byly nalezeny a~opraveny tyto chyby:

\begin{enumerate}
    \item \textbf{\texttt{single-galerie.php}} -- sezóna se načítala z~meta pole místo z~taxonomie. Opraveno na čtení z~taxonomie \texttt{sezona}.
    \item \textbf{\texttt{front-page.php}} -- Bootstrap třída \texttt{.card} přidávala nechtěný rámeček. Opraveno přidáním CSS \texttt{border: none}.
    \item \textbf{\texttt{front-page.php}} -- horizontální scrollbar na mobilních zařízeních. Opraveno přidáním Bootstrap třídy \texttt{g-0} pro odstranění gutterů.
\end{enumerate}

% ========================================================================
\clearpage
\section{Příručka administrátora}

\subsection{Návod pro správce -- administrace WordPress}

\subsubsection{Přihlášení}

Otevřete přihlašovací stránku webu a~přihlaste se. Výchozí adresa WordPressu je \texttt{/wp-login.php}, ale na produkčním serveru je pomocí pluginu \textbf{WPS Hide Login} změněna na vlastní slug (viz sekce~3.12). V~levém menu uvidíte sekce: \textbf{Zápasy}, \textbf{Týmy}, \textbf{Hráči}, \textbf{Galerie}, \textbf{Sponzoři}, \textbf{Kontakty}.

\subsubsection{Správa zápasů}

\begin{enumerate}
    \item Klikněte \textbf{Zápasy $\rightarrow$ Přidat zápas}
    \item Vyplňte titulek (např.\ \uv{TJ Slavoj Mýto vs SK Nepomuk})
    \item V~sekci \textbf{Detail zápasu} vyplňte datum, čas, domácí tým, hosty, skóre a~střelce
    \item V~pravém panelu přiřaďte \textbf{kategorii týmu} a~\textbf{sezónu}
    \item Klikněte \textbf{Publikovat}
\end{enumerate}

Pro zadání výsledku odehraného zápasu: otevřete existující zápas, vyplňte \textbf{Skóre} a~\textbf{Střelce}, klikněte \textbf{Aktualizovat}.

\subsubsection{Správa týmů}

\begin{enumerate}
    \item Klikněte \textbf{Týmy $\rightarrow$ Přidat tým}
    \item Vyplňte název, identifikátor týmu (slug, např.\ \texttt{muzi-a}), trenéry
    \item Nastavte \textbf{Obrázek příspěvku} (týmová fotografie)
    \item Přiřaďte sezónu a~kategorii
    \item \textbf{Počet hráčů} se vypočítá automaticky z~přiřazených hráčů
\end{enumerate}

\subsubsection{Správa hráčů}

\begin{enumerate}
    \item Klikněte \textbf{Hráči $\rightarrow$ Přidat hráče}
    \item Vyplňte jméno, číslo dresu (1--99), rok narození, slug týmu
    \item Přiřaďte pozici (Brankář / Hráč v~poli) a~sezónu
\end{enumerate}

\subsubsection{Správa galerií}

\begin{enumerate}
    \item Klikněte \textbf{Galerie $\rightarrow$ Přidat album}
    \item Vyplňte název alba
    \item V~obsahu příspěvku klikněte \textbf{Přidat média $\rightarrow$ Vytvořit galerii} $\rightarrow$ vyberte fotky $\rightarrow$ \textbf{Vložit galerii}
    \item Nastavte \textbf{Obrázek příspěvku} jako náhled alba
    \item V~pravém panelu přiřaďte \textbf{Kategorii týmu}
\end{enumerate}

\subsubsection{Správa sponzorů}

\begin{enumerate}
    \item Klikněte \textbf{Sponzoři $\rightarrow$ Přidat sponzora}
    \item Vyplňte název firmy / partnera jako titulek příspěvku
    \item V~sekci \textbf{Detail sponzora} zadejte \textbf{Webové stránky} (URL, např.\ \texttt{https://firma.cz})
    \item Nastavte \textbf{Obrázek příspěvku} -- logo sponzora
    \item Klikněte \textbf{Publikovat}
\end{enumerate}

\subsubsection{Správa kontaktů}

\begin{enumerate}
    \item Klikněte \textbf{Kontakty $\rightarrow$ Přidat kontakt}
    \item Vyplňte jméno osoby jako titulek příspěvku (např.\ \uv{Jaroslav Bejček})
    \item Vyplňte pole v~sekci \textbf{Detail kontaktu}:

    \smallskip
    \begin{tabularx}{\linewidth}{l X l}
    \toprule
    \textbf{Pole} & \textbf{Popis} & \textbf{Příklad} \\
    \midrule
    Funkce / Pozice & Role v~klubu & \texttt{Předseda klubu} \\
    Telefon & Kontaktní telefon & \texttt{+420 777 000 000} \\
    E-mail & Kontaktní e-mail & \texttt{predseda@slavojmyto.cz} \\
    Pořadí zobrazení & Číslo pro řazení (0 = první) & \texttt{1} \\
    \bottomrule
    \end{tabularx}
    \smallskip

    \item Klikněte \textbf{Publikovat}
\end{enumerate}

\begin{quote}
\textbf{Tip:} Nižší číslo v~poli \uv{Pořadí zobrazení} = kontakt se zobrazí výše v~seznamu na webu.
\end{quote}

\subsubsection{Správa aktualit}

Aktuality využívají nativní WordPress příspěvky a~zobrazují se na úvodní stránce.

\begin{enumerate}
    \item Klikněte \textbf{Příspěvky $\rightarrow$ Přidat příspěvek}
    \item Napište titulek a~text aktuality
    \item Nastavte \textbf{Obrázek příspěvku} (náhled na homepage)
    \item Klikněte \textbf{Publikovat}
\end{enumerate}

\subsection{Příručka pro návštěvníka webu}

Web TJ Slavoj Mýto je rozdělen do přehledných sekcí přístupných z~hlavního menu. Níže je stručný popis uživatelského rozhraní.

\textbf{Navigace:} Hlavní menu v~záhlaví stránky obsahuje odkazy na Domů, Zápasy, Týmy, Galerie, Sponzoři, Kontakty a~Historie. Na mobilních zařízeních se menu sbalí do hamburger ikony, kliknutím se rozbalí.

\textbf{Úvodní stránka (\texttt{/}):} Zobrazuje nejbližší nadcházející zápas, poslední výsledky a~nejnovější aktuality. Slouží jako rozcestník k~ostatním sekcím webu.

\textbf{Zápasy (\texttt{/zapasy/}):} Přehled zápasů s~filtry podle týmu, sezóny a~stavu. Každý zápas zobrazuje datum, týmy, skóre (barevně -- zelená = výhra, červená = prohra) a~střelce. Kliknutím na zápas se zobrazí jeho detail.

\textbf{Týmy (\texttt{/tymy/}):} Karty týmů s~fotografiemi. Kliknutím na tým se zobrazí detail se soupiskou, realizačním týmem a~nadcházejícími zápasy.

\textbf{Galerie (\texttt{/galerie/}):} Přehled fotoalb s~náhledovými obrázky. Kliknutím na album se zobrazí fotografie s~lightboxem (zvětšení fotografie na celou obrazovku).

\textbf{Sponzoři (\texttt{/sponzori/}):} Přehled partnerů klubu s~logy. Kliknutím na logo se otevře web sponzora.

\textbf{Kontakty (\texttt{/kontakty/}):} Seznam členů vedení klubu s~funkcí, telefonem a~e-mailem.

\textbf{Historie (\texttt{/historie/}):} Stránka s~historií a~tradicemi klubu.

% ========================================================================
\clearpage
\section{Nasazení}

Tato kapitola popisuje nasazení webu na produkční hosting. Doporučeným řešením je český hosting \textbf{Wedos} s~podporou PHP, MariaDB a~FTP přístupu.

\subsection{Nasazení na hosting Wedos}

\textbf{Požadavky na hosting:}

\begin{table}[H]
\centering
\begin{tabular}{ll}
\toprule
\textbf{Požadavek} & \textbf{Minimální hodnota} \\
\midrule
PHP & $\geq$ 8.0 \\
MariaDB & $\geq$ 10.4 \\
Diskový prostor & $\geq$ 1\,GB \\
FTP přístup & Ano \\
\bottomrule
\end{tabular}
\caption{Požadavky na hosting}
\end{table}

\textbf{Postup nasazení:}

\begin{enumerate}
    \item \textbf{Vytvoření databáze} -- V~klientské zóně Wedos (zákaznická administrace) vytvořte novou MariaDB databázi. Poznamenejte si název databáze, uživatelské jméno, heslo a~hostname (obvykle \texttt{localhost} nebo adresa uvedená v~administraci Wedos).

    \item \textbf{Příprava WordPressu} -- Stáhněte WordPress z~\url{https://wordpress.org/download/} a~rozbalte na lokálním počítači. Do složky WordPressu zkopírujte téma a~plugin:

\begin{lstlisting}
web/wp-content/themes/tj-slavoj-myto/
  -> wp-content/themes/tj-slavoj-myto/
web/wp-content/plugins/slavoj-custom-fields/
  -> wp-content/plugins/slavoj-custom-fields/
\end{lstlisting}

    \item \textbf{Nastavení \texttt{wp-config.php}} -- Přejmenujte \texttt{wp-config-sample.php} na \texttt{wp-config.php} a~vyplňte údaje databáze:

\begin{lstlisting}[language=PHP]
define( 'DB_NAME',     'nazev_databaze' );
define( 'DB_USER',     'uzivatel_databaze' );
define( 'DB_PASSWORD', 'heslo_databaze' );
define( 'DB_HOST',     'localhost' );
define( 'DB_CHARSET',  'utf8mb4' );
\end{lstlisting}

    Vygenerujte bezpečnostní klíče na \url{https://api.wordpress.org/secret-key/1.1/salt/} a~vložte je do souboru.

    \item \textbf{Nahrání souborů přes FTP} -- Připojte se k~FTP serveru Wedos pomocí klienta (např.\ FileZilla). FTP údaje najdete v~klientské zóně Wedos. Nahrajte \textbf{obsah} složky WordPressu (všechny soubory a~podsložky) do kořenového adresáře webu (\texttt{www/}).

    \item \textbf{Instalace WordPressu} -- Otevřete v~prohlížeči adresu vašeho webu (např.\ \texttt{https://vasedomena.cz/}). WordPress spustí instalačního průvodce -- vyplňte název webu (\texttt{TJ Slavoj Mýto}), administrátorský účet a~e-mail.

    \item \textbf{Aktivace tématu a~pluginu} -- V~administraci aktivujte téma \textbf{TJ Slavoj Mýto} (Vzhled $\rightarrow$ Témata) a~plugin \textbf{Slavoj Custom Fields} (Pluginy). Nastavte trvalé odkazy na \textbf{Název příspěvku} (Nastavení $\rightarrow$ Trvalé odkazy).

    \item \textbf{Prvotní nastavení obsahu} -- Přejděte na \textbf{Nástroje $\rightarrow$ Slavoj nastavení} a~klikněte \uv{Vytvořit výchozí hodnoty taxonomií} a~\uv{Vytvořit stránky webu}. Nastavte hlavní stránku (Nastavení $\rightarrow$ Čtení $\rightarrow$ Statická stránka $\rightarrow$ Domů) a~vytvořte hlavní menu.

    \item \textbf{SSL certifikát} -- V~klientské zóně Wedos aktivujte SSL certifikát (Let's Encrypt -- zdarma). V~WordPress admin změňte URL webu na \texttt{https://} (Nastavení $\rightarrow$ Obecné).
\end{enumerate}

\subsection{Další možnosti nasazení}

Kromě Wedos lze web nasadit na jakýkoli hosting s~podporou PHP $\geq$~8.0 a~MySQL/MariaDB (např.\ Active24, Blueboard). Postup je shodný -- liší se pouze přístup ke klientské zóně a~konkrétní nastavení FTP a~databáze. Pro lokální vývoj je doporučen XAMPP~\cite{xampp}.

% ========================================================================
\clearpage
\section{Závěr}

Cílem této maturitní práce bylo vytvořit moderní webovou prezentaci pro fotbalový klub TJ Slavoj Mýto, která umožní správci klubu spravovat obsah bez znalosti programování. Tento cíl byl splněn.

Výsledný web je postaven na systému WordPress s~vlastním tématem \texttt{tj-slavoj-myto} a~pluginem \texttt{slavoj-custom-fields}. Místo závislosti na externích pluginech byla klíčová funkcionalita implementována vlastním kódem -- registrace šesti typů obsahu, čtyř taxonomií, administrační meta boxy, vlastní sloupce, filtry a~uživatelská role \uv{Správce obsahu}.

Web je plně responzivní díky frameworku Bootstrap~5 a~mobile-first přístupu ke stylování. Kód dodržuje bezpečnostní standardy WordPressu (sanitizace, escapování, nonce) a~minimalizuje použití JavaScriptu.

\subsection*{Splnění bodů zadání}

\begin{table}[H]
\centering
\begin{tabularx}{\textwidth}{X l}
\toprule
\textbf{Bod zadání} & \textbf{Stav} \\
\midrule
Figma design (desktop + mobilní) & Vytvořen \\
WordPress téma s~Bootstrap~5 & Implementováno \\
Custom post types (zápasy, týmy, hráči\ldots) & 6~CPT + 4~taxonomie \\
Filtrování podle sezóny a~týmu & GET filtry na archivních stránkách \\
Galerie fotografií & CPT s~nativní WP galerií + lightbox \\
Optimalizace výkonu & Lazy loading, CDN, minimální JS \\
Uživatelské role & Administrátor + Správce obsahu \\
\bottomrule
\end{tabularx}
\caption{Splnění bodů zadání}
\end{table}

\subsection*{Návrhy na rozšíření}

\begin{enumerate}
    \item \textbf{Kalendářní modul} -- vizuální kalendář zápasů (doporučen plugin The Events Calendar~\cite{eventscalendar})
    \item \textbf{SEO optimalizace} -- integrace pluginu Rank Math~\cite{rankmath} nebo Yoast SEO~\cite{yoast}
    \item \textbf{Cache a~komprese} -- nasazení WP Super Cache~\cite{wpsupercache} a~Smush~\cite{smush} pro produkční prostředí
    \item \textbf{Vícejazyčnost} -- příprava pro anglickou verzi webu
    \item \textbf{Statistiky hráčů} -- rozšíření dat hráčů o~góly, asistence a~zápasy
\end{enumerate}

% ========================================================================
% SEZNAM POUŽITÉ LITERATURY
% ========================================================================

\newpage
\phantomsection
\addcontentsline{toc}{section}{Seznam použité literatury}
\begin{thebibliography}{99}

\bibitem{zadani}
BEJČEK, Lukáš. \textit{Zadání maturitní práce: Webová prezentace fotbalového klubu s~databází zápasů a~správou obsahu ve~WordPressu} [dokument]. 2025. Soubor \texttt{ZADANI-MATURITNI-PRACE.md} v~kořeni repozitáře.

\bibitem{wordpress}
WORDPRESS FOUNDATION. \textit{WordPress} [software]. Verze 6.x. San Francisco (CA): WordPress Foundation, 2024 [cit. 2026-03-27]. Dostupné z: \url{https://wordpress.org}

\bibitem{wpdev}
WORDPRESS FOUNDATION. \textit{WordPress Developer Resources} [online]. San Francisco (CA): WordPress Foundation, 2024 [cit. 2026-03-27]. Dostupné z: \url{https://developer.wordpress.org}

\bibitem{bootstrap}
OTTO, Mark a~Jacob THORNTON. \textit{Bootstrap} [software]. Verze 5.3.3. Bootstrap Team, 2024 [cit. 2026-03-27]. Dostupné z: \url{https://getbootstrap.com}

\bibitem{jsdelivr}
JSDELIVR. \textit{jsDelivr -- free CDN for open source} [online]. 2024 [cit. 2026-03-27]. Dostupné z: \url{https://www.jsdelivr.com}

\bibitem{xampp}
APACHE FRIENDS. \textit{XAMPP} [software]. Apache Friends, 2024 [cit. 2026-03-27]. Dostupné z: \url{https://www.apachefriends.org}

\bibitem{figma}
FIGMA, INC. \textit{Figma} [software]. San Francisco (CA): Figma, 2024 [cit. 2026-03-27]. Dostupné z: \url{https://www.figma.com}

\bibitem{cptui}
FLAVOR, WebDevStudios. Custom Post Type UI. In: \textit{WordPress Plugin Directory} [online]. San Francisco (CA): WordPress Foundation, 2024 [cit. 2026-03-27]. Dostupné z: \url{https://wordpress.org/plugins/custom-post-type-ui/}

\bibitem{acf}
WP ENGINE. Advanced Custom Fields. In: \textit{WordPress Plugin Directory} [online]. San Francisco (CA): WordPress Foundation, 2024 [cit. 2026-03-27]. Dostupné z: \url{https://wordpress.org/plugins/advanced-custom-fields/}

\bibitem{facetwp}
FACETWP, LLC. \textit{FacetWP -- Advanced Filtering for WordPress} [online]. 2024 [cit. 2026-03-27]. Dostupné z: \url{https://facetwp.com}

\bibitem{roleeditor}
FLAVOR, Vladimir Garagulja. User Role Editor. In: \textit{WordPress Plugin Directory} [online]. San Francisco (CA): WordPress Foundation, 2024 [cit. 2026-03-27]. Dostupné z: \url{https://wordpress.org/plugins/user-role-editor/}

\bibitem{lazyload}
WORLD WIDE WEB CONSORTIUM (W3C). \textit{HTML Living Standard: Lazy loading} [online]. 2024 [cit. 2026-03-27]. Dostupné z: \url{https://html.spec.whatwg.org/multipage/urls-and-fetching.html}

\bibitem{themehandbook}
WORDPRESS FOUNDATION. \textit{Theme Developer Handbook} [online]. San Francisco (CA): WordPress Foundation, 2024 [cit. 2026-03-27]. Dostupné z: \url{https://developer.wordpress.org/themes/}

\bibitem{jsreport}
BEJČEK, Lukáš. \textit{JavaScript v~projektu TJ Slavoj Mýto} [dokument]. 2026. Soubor \texttt{docs/pracovni/12-javascript.md} v~repozitáři.

\bibitem{git}
SOFTWARE FREEDOM CONSERVANCY. \textit{Git} [software]. Verze 2.x. Software Freedom Conservancy, 2024 [cit. 2026-03-27]. Dostupné z: \url{https://git-scm.com}

\bibitem{docker}
DOCKER, INC. \textit{Docker} [software]. San Francisco (CA): Docker, 2024 [cit. 2026-03-27]. Dostupné z: \url{https://www.docker.com}

\bibitem{eventscalendar}
THE EVENTS CALENDAR. The Events Calendar. In: \textit{WordPress Plugin Directory} [online]. San Francisco (CA): WordPress Foundation, 2024 [cit. 2026-03-27]. Dostupné z: \url{https://wordpress.org/plugins/the-events-calendar/}

\bibitem{rankmath}
FLAVOR, Starter Sites. Rank Math SEO. In: \textit{WordPress Plugin Directory} [online]. San Francisco (CA): WordPress Foundation, 2024 [cit. 2026-03-27]. Dostupné z: \url{https://wordpress.org/plugins/seo-by-rank-math/}

\bibitem{yoast}
YOAST BV. Yoast SEO. In: \textit{WordPress Plugin Directory} [online]. San Francisco (CA): WordPress Foundation, 2024 [cit. 2026-03-27]. Dostupné z: \url{https://wordpress.org/plugins/wordpress-seo/}

\bibitem{wpsupercache}
FLAVOR, Developer. WP Super Cache. In: \textit{WordPress Plugin Directory} [online]. San Francisco (CA): WordPress Foundation, 2024 [cit. 2026-03-27]. Dostupné z: \url{https://wordpress.org/plugins/wp-super-cache/}

\bibitem{smush}
FLAVOR, Developer. Smush. In: \textit{WordPress Plugin Directory} [online]. San Francisco (CA): WordPress Foundation, 2024 [cit. 2026-03-27]. Dostupné z: \url{https://wordpress.org/plugins/wp-smushit/}

\end{thebibliography}

% ========================================================================
% PŘÍLOHY
% ========================================================================

\newpage
\appendix

\section{Kopie zadání práce}

\textbf{Jméno a~příjmení:} Lukáš Bejček\\
\textbf{Studijní obor:} Informační technologie\\
\textbf{Název tématu:} Webová prezentace fotbalového klubu s~databází zápasů a~správou obsahu ve~WordPressu

\subsection*{Cíl práce}

Cílem práce je vytvořit moderní informační web pro fotbalový klub, který bude přehledně zobrazovat zápasy, týmy, historii, partnery a~aktuální informace o~dění v~klubu. Web bude postaven na vlastním grafickém návrhu vytvořeném ve Figmě a~nasazeném na systému WordPress s~využitím vlastní šablony založené na Bootstrapu~5. Web bude připraven pro správu obsahu vedením klubu nebo pověřeným administrátorem.

\subsection*{Zásady pro vypracování}

\begin{enumerate}
    \item Vytvoření grafického návrhu ve Figmě. Návrh musí obsahovat samostatné varianty pro desktop i~mobilní zařízení, včetně hlavní stránky i~podstránek. Grafický návrh musí zahrnovat také vizuální podobu použitých pluginů a~komponent, aby zapadaly do celkového designu stránky.

    \item Tvorba šablony v~Bootstrapu~5 a~nasazení ve WordPressu. Použití souborů \texttt{header.php}, \texttt{footer.php}, \texttt{single.php}, \texttt{archive.php}, \texttt{functions.php}, \texttt{styles.css} atd.

    \item Zápasy a~týmy jako vlastní typy obsahu
    \begin{enumerate}[label=\alph*)]
        \item Vytvoření typu příspěvku \uv{Zápas} s~poli: datum, čas, tým, soupeř, výsledek, střelci, sezóna, status zápasu
        \item Vytvoření typu příspěvku \uv{Tým} s~poli: název, sezóna, trenér, asistenti, seznam hráčů
        \item Pluginy: Custom Post Type UI, Advanced Custom Fields, Pods (volitelné)
    \end{enumerate}

    \item Filtrování a~výpis podle sezóny nebo týmů
    \begin{enumerate}[label=\alph*)]
        \item Možnost filtrování zápasů a~týmů podle ročníku, stavu, nebo kategorie
        \item Pluginy: FacetWP, Search \& Filter Pro
    \end{enumerate}

    \item Kalendářový modul
    \begin{enumerate}[label=\alph*)]
        \item Implementace kalendáře, který umožní vizuální zobrazení nadcházejících i~odehraných zápasů
        \item Možnost filtrování podle týmů, kategorií nebo sezón
        \item Pluginy: The Events Calendar, Event Organiser, Simple Calendar
    \end{enumerate}

    \item Galerie
    \begin{enumerate}[label=\alph*)]
        \item Vytvoření galerie s~možností seskupení podle názvu nebo roku
        \item Pluginy: Modula, Envira Gallery, NextGEN Gallery
    \end{enumerate}

    \item Optimalizace výkonu a~načítání
    \begin{enumerate}[label=\alph*)]
        \item Cache (mezipaměť) -- ukládání stránek do mezipaměti pro rychlejší načítání
        \item Lazy loading -- odložené načítání obrázků a~externího obsahu až při posunu stránky
        \item Minifikace -- zmenšení velikosti CSS a~JS souborů
        \item Optimalizace obrázků -- komprese a~správný formát (např.\ WebP)
        \item Kombinace CSS/JS -- snížení počtu HTTP požadavků
        \item Pluginy: WP Super Cache, LiteSpeed Cache, Smush
    \end{enumerate}

    \item SEO a~analytika
    \begin{enumerate}[label=\alph*)]
        \item Pluginy: Yoast SEO, Rank Math, propojení s~Google Analytics
    \end{enumerate}

    \item Hosting a~zabezpečení
    \begin{enumerate}[label=\alph*)]
        \item Hosting, HTTPS certifikát, ochrana proti útokům a~zálohování
    \end{enumerate}

    \item Uživatelské role a~oprávnění
    \begin{enumerate}[label=\alph*)]
        \item Administrátor, přispěvatel -- různá oprávnění pro správu obsahu
        \item Pluginy: User Role Editor, Members
    \end{enumerate}

    \item Odevzdávané výstupy
    \begin{enumerate}[label=\alph*)]
        \item Podrobná dokumentace všech použitých pluginů
        \item Zdrojový kód WordPress šablony a~popis všech jejích úprav
        \item Funkční web (na hostingu nebo jako lokalizovaný projekt s~návodem ke spuštění a~veškerými přístupovými údaji)
    \end{enumerate}
\end{enumerate}

\textbf{Rozsah dokumentace:} 15--50 stran\\
\textbf{Forma MP:} elektronická

% ========================================================================
\section{Přístupové údaje}

\begin{quote}
\textbf{Upozornění:} Tyto údaje jsou určeny výhradně pro lokální testování. Při nasazení na produkční server je nutné zvolit silná a~unikátní hesla.
\end{quote}

\subsection*{WordPress administrace}

\textbf{URL:} \texttt{http://localhost/fotbal\_club/wp-admin/}

\begin{table}[H]
\centering
\small
\begin{tabularx}{\textwidth}{l l l X}
\toprule
\textbf{Role} & \textbf{Uživatel} & \textbf{Heslo} & \textbf{Oprávnění} \\
\midrule
Administrátor & \texttt{admin} & \texttt{admin123} & Plný přístup k~celému WordPressu \\
Správce obsahu & \texttt{redaktor} & \texttt{redaktor123} & Editace CPT, nahrávání médií, bez nastavení WP a~pluginů \\
\bottomrule
\end{tabularx}
\caption{Přístupové údaje WordPress}
\end{table}

\textbf{Administrátor} má plný přístup ke všem funkcím WordPressu včetně nastavení, pluginů, témat a~všech typů obsahu.

\textbf{Správce obsahu} může:
\begin{itemize}
    \item Přidávat, editovat a~mazat zápasy, týmy, hráče, galerie, sponzory a~kontakty
    \item Nahrávat obrázky a~média
    \item Spravovat kategorie a~taxonomie
\end{itemize}

Správce obsahu \textbf{nemůže}:
\begin{itemize}
    \item Měnit nastavení WordPressu
    \item Instalovat nebo mazat pluginy a~témata
    \item Spravovat uživatelské účty
\end{itemize}

\subsection*{Databáze (MySQL -- XAMPP)}

\begin{table}[H]
\centering
\begin{tabular}{ll}
\toprule
\textbf{Parametr} & \textbf{Hodnota} \\
\midrule
Server & \texttt{localhost} \\
Port & \texttt{3306} \\
Uživatel & \texttt{root} \\
Heslo & \textit{(prázdné -- výchozí XAMPP)} \\
Databáze & \texttt{slavoj\_myto} \\
phpMyAdmin & \texttt{http://localhost/phpmyadmin/} \\
\bottomrule
\end{tabular}
\caption{Přístupové údaje databáze}
\end{table}

\end{document}
